% -----------------------------------------------------------------------------------
\section{Row-baseline invariance and contraction}\label{sec:translations}

Adding row-translation invariance changes the picture. It produces a genuine impossibility among
four independent axioms, none of which concerns Pareto, ordering, or ex-ante value.

\subsection{The circulant theorem}

The proof embeds an arbitrary matrix in a matrix with a transitive symmetry group. It rests on a
distinct-differences lemma of the kind used by \textcite{finkruiz1986} to embed graphs in
circulants.

\begin{lemma}[Distinct differences]\label{lem:sidon}
Let $k\ge2$, $v_i=2^i$ for $i=0,\dots,k-1$, and $m>2(2^{k-1}-1)$. The $k(k-1)$ ordered differences
$v_j-v_i$, $i\ne j$, are distinct and nonzero modulo $m$.
\end{lemma}

\begin{proof}
For $j>i$, $v_j-v_i=2^i(2^{j-i}-1)$; the power of two and the odd part determine $i$ and $j-i$, so the
positive differences are distinct integers, and so are their negatives. Every difference has
absolute value at most $D=2^{k-1}-1$, so two distinct differences differ by a nonzero integer of
absolute value at most $2D<m$, and no difference is divisible by $m$.
\end{proof}

\begin{lemma}[Circulant embedding]\label{lem:embed}
Every matrix $Z\in\Rl^{k\times k}$ with zero diagonal is a principal submatrix of a circulant matrix
$A_{uv}=f(v-u\bmod m)$ on $\mathbb Z/m\mathbb Z$ with $f(0)=0$.
\end{lemma}

\begin{proof}
With $v_i$ and $m$ as in Lemma~\ref{lem:sidon}, set $f(0)=0$ and $f(v_j-v_i)=Z_{ij}$ for $i\ne j$, which
is well defined by the lemma, and give $f$ arbitrary values on the unused residues. Then
$A_{v_iv_j}=f(v_j-v_i)=Z_{ij}$.
\end{proof}

\begin{theorem}[Neutrality, row translations, and contraction force universal choice]\label{thm:circulant}
On the domain of all finite real square matrices, every rule satisfying \ax{N}, \ax{R}, and \ax{C}
(principal restriction) has $G(M)=I$ for every matrix $(I,M)$. Consequently
$\{\ax{N},\ax{R},\ax{C},\ax{Q}\}$ is inconsistent.
\end{theorem}

\begin{proof}
Let $(I,M)$ be given with $|I|=k\ge2$ (the case $k=1$ is trivial), and let $Z=Z(M)$, which has zero
diagonal; by \ax{R}, $G(Z)=G(M)$. By Lemma~\ref{lem:embed} and \ax{N} (relabelling the vertices of the
circulant so that the embedded ones carry the labels of $I$), $Z$ is the principal restriction to $I$
of a matrix $A$ whose symmetry group acts transitively on its indices: each cyclic shift
$u\mapsto u+1$ of $\mathbb Z/m\mathbb Z$, transported by the relabelling, maps $A$ to itself. By \ax{N},
$G(A)$ is invariant under these symmetries, and since it is nonempty it contains every index. By
\ax{C}, $G(Z)=G(A|_I)\supseteq G(A)\cap I=I$. Hence $G(M)=I$, and no two matrices on the same index set
have different choice sets.
\end{proof}

Without \ax{O} the conclusion is universal \emph{choice}, not universal indifference: incomparabilities
may remain. The symmetry step is standard in the theory of tournament solutions, where neutral,
nonempty solutions must select every vertex of a vertex-transitive tournament
\autocite{brandt2011,brandtbrillharrenstein2016}; the embedding has the precedent noted above. What
the theorem adds is the use of contraction to carry universal choice from symmetric matrices to all
matrices, and of row-baseline invariance to reach every matrix from a zero-diagonal one. No
published statement of this conclusion was found in the literature searched; its ingredients are
established, and no priority is claimed.

\begin{proposition}[Independence]\label{prop:indep}
Each three-element subset of $\{\ax{N},\ax{R},\ax{C},\ax{Q}\}$ is satisfiable, so the family is
minimally inconsistent:
\begin{center}
\small
\begin{tabular}{@{}ll@{}}
  \toprule
  omitted axiom & rule satisfying the other three \\
  \midrule
  \ax{N} & responsive fixed-label rule \\
  \ax{R} & diagonal ranking \\
  \ax{C} & row-range ranking, $r_i=\max_jM_{ij}-\min_jM_{ij}$ \\
  \ax{Q} & universal indifference \\
  \bottomrule
\end{tabular}
\end{center}
All four rules also satisfy \ax{F} and \ax{O}.
\end{proposition}

\begin{proof}
Row ranges are unchanged by row translations, multiplied by $a$ under \ax{F}, and permuted under
relabelling; the rule is nonconstant. Diagonal ranking satisfies \ax{C} by Lemma~\ref{lem:bimplies}.
The \emph{responsive fixed-label rule} fixes two labels $u,v$ and a well-order of $\mathbb L$ in which
$u,v$ form an adjacent block above all other labels. When both occur, $u$ and $v$ are compared by the
sign of $M_{uu}-M_{uv}$ (a tie if it is zero); all other comparisons follow the fixed order. The
comparison of $u$ and $v$ uses a within-row difference of their principal block, so the rule satisfies
\ax{R}, \ax{F}, and \ax{B}, hence \ax{C}; it is responsive, so \ax{Q} holds, but it is not neutral.
\end{proof}

If neutrality is regarded as part of the reduction \ax{A0} rather than as an axiom, the same result is
a minimal three-axiom family $\{\ax{R},\ax{C},\ax{Q}\}$ relative to that background.

\begin{proposition}\label{prop:NRBQ}
$\{\ax{N},\ax{R},\ax{B},\ax{Q}\}$ is minimally inconsistent, and so is $\{\ax{R},\ax{D}\}$.
\end{proposition}

\begin{proof}
Suppose $i\succ_Mj$ in some matrix. By \ax{B} the same holds in $M|_{\{i,j\}}$, and by \ax{R} in its
row normalisation $H=\bigl(\begin{smallmatrix}0&x\\y&0\end{smallmatrix}\bigr)$, with $i$ first. Let
\[
  C=\begin{pmatrix}0&x&y\\y&0&x\\x&y&0\end{pmatrix}.
\]
Its principal blocks on $(1,2)$, $(2,3)$, and $(3,1)$ all equal $H$, so \ax{N} and \ax{B} give
$1\succ2$, $2\succ3$, and $3\succ1$, and $G(C)=\varnothing$. Hence no strict comparison ever occurs and
$G(M)=I$ always, contradicting \ax{Q}. Minimality follows from the witnesses of
Proposition~\ref{prop:indep}, with row-range ranking now serving as the rule that omits \ax{B}. For $\{\ax{R},\ax{D}\}$: adding $2$ to the
second row of $\bigl(\begin{smallmatrix}1&0\\0&0\end{smallmatrix}\bigr)$ changes the unique diagonal
maximiser from option $1$ to option $2$, while \ax{R} keeps the choice fixed.
\end{proof}

Unlike Theorem~\ref{thm:circulant}, this argument uses the link between comparisons and nonempty
maximal choice fixed in Definition~\ref{def:rule}.

\begin{theorem}[Map with row translations]\label{thm:augmented}
Adding \ax{R} to the family of Theorem~\ref{thm:pairwise} introduces exactly three new minimal
inconsistent subsets,
\[
  \{\ax{R},\ax{D}\},\qquad\{\ax{N},\ax{R},\ax{B},\ax{Q}\},\qquad\{\ax{N},\ax{R},\ax{C},\ax{Q}\},
\]
and the maximal consistent subsets become
\begin{center}
\small
\begin{tabular}{@{}ll@{}}
  \toprule
  witness rule & maximal consistent subset \\
  \midrule
  column sums & $\{\ax{N},\ax{F},\ax{P},\ax{R},\ax{O},\ax{Q}\}$ \\
  diagonal ranking & $\{\ax{N},\ax{F},\ax{D},\ax{B},\ax{S},\ax{C},\ax{O},\ax{Q}\}$ \\
  row ranges & $\{\ax{N},\ax{F},\ax{R},\ax{S},\ax{O},\ax{Q}\}$ \\
  responsive fixed-label rule & $\{\ax{F},\ax{R},\ax{B},\ax{S},\ax{C},\ax{O},\ax{Q}\}$ \\
  universal indifference & $\{\ax{N},\ax{F},\ax{R},\ax{B},\ax{S},\ax{C},\ax{O}\}$ \\
  \bottomrule
\end{tabular}
\end{center}
\end{theorem}

\begin{proof}
Column sums satisfy \ax{R} because a row translation shifts every column sum by $\sum_kb_k$; row ranges
satisfy \ax{S} because each score depends only on its own row. Let a set avoid all minimal
inconsistent subsets. If it contains \ax{P}, its pairwise conflicts confine it to the first row. If it
omits \ax{P} and \ax{R}, it fits the second. If it contains \ax{R}, it omits \ax{D}; if it contains both
\ax{N} and \ax{Q}, the two four-element obstructions exclude \ax{B} and \ax{C}, leaving the third row;
without \ax{N} it fits the fourth, and without \ax{Q} the fifth.
\end{proof}

\subsection{An exact domain condition}

Theorem~\ref{thm:circulant} needs matrices of unbounded size: with a fixed maximum size $n$, the rule
that chooses everything below size $n$ and maximises row ranges at size $n$ satisfies \ax{N}, \ax{R},
\ax{C}, and \ax{Q}. The exact requirement on the domain can be stated without quantifying over
rules. Let $\mathcal D$ be a set of matrices closed under relabelling, and interpret each invariance as
applying whenever both the matrix and its transform belong to $\mathcal D$. Form a directed graph on the
\emph{pointed matrices} $\mathcal P_{\mathcal D}=\{(M,i):M\in\mathcal D,\ i\in I_M\}$, with
\begin{enumerate}[label=(\roman*)]
  \item edges in both directions between $(M,i)$ and $(\sigma\cdot M,\sigma(i))$ for admissible relabellings;
  \item edges in both directions between $(M,i)$ and $(M+b\one^\top,i)$ for admissible row translations;
  \item extension edges $(M|_K,i)\to(M,i)$ for $i\in K$, whenever $M|_K$ and $M$ are in $\mathcal D$.
\end{enumerate}
Let $\Reach(x)$ be the set of pointed matrices reachable from $x$ (including $x$), and let
$\Phi_M=\{(M,i):i\in I_M\}$ be the fibre of $M$.

\begin{theorem}[Exact domain criterion]\label{thm:DR}
On $\mathcal D$, the axioms $\ax{N},\ax{R},\ax{C},\ax{Q}$ are jointly inconsistent if and only if
\[
  \text{for every }x\in\mathcal P_{\mathcal D}\text{ there is }M\in\mathcal D\text{ with }\Phi_M\subseteq\Reach(x).
  \tag{DR}
\]
\end{theorem}

\begin{proof}
Say $(M,i)$ is \emph{rejected} if $i\notin G(M)$. Rejection propagates along relabelling and
translation edges by \ax{N} and \ax{R}, and along extension edges by the contrapositive of \ax{C}. Under
(DR), a single rejection at $x$ would reject the entire fibre of some $M$, contradicting $G(M)\neq\varnothing$;
so nothing is ever rejected, and \ax{Q} fails. Conversely, if (DR) fails at $x=(M_0,i_0)$, set
\[
  G_x(M)=\{i:(M,i)\notin\Reach(x)\}.
\]
Each $G_x(M)$ is nonempty because no fibre lies in $\Reach(x)$. Since relabelling and translation
edges are bidirectional, $\Reach(x)$ is closed under them, which gives \ax{N} and \ax{R}. If
$i\in G_x(M)\cap K$ but $i\notin G_x(M|_K)$, the extension edge would put $(M,i)$ in $\Reach(x)$; so \ax{C}
holds. Finally $i_0\notin G_x(M_0)$, so $G_x(M_0)$ is a nonempty proper subset of $I_{M_0}$, and some
relabelling of $M_0$ on the same index set moves it, which gives \ax{Q}. Tying chosen options above
rejected ones yields \ax{O} as well.
\end{proof}

The circulant condition is sufficient for (DR) but not necessary. If \ax{F} is also imposed, admissible
affine transformations are added as bidirectional edges.

\subsection{Physical restrictions on problems}

Does the richness that (DR) demands survive natural restrictions on the problems themselves?

\begin{proposition}[Bounded payoffs]\label{prop:bounded}
On the domain of $[0,1]$-valued matrices, $\{\ax{N},\ax{R},\ax{C},\ax{Q}\}$ is consistent; adding \ax{F}
makes it inconsistent again.
\end{proposition}

\begin{proof}
Fix $r\in[\tfrac12,1)$ and put $i\to j$ if $M_{ii}-M_{ij}>r$ and $M_{ji}-M_{jj}>r$. An edge forces
$M_{ii}>r\ge\tfrac12\ge1-r>M_{jj}$, so no vertex has both an incoming and an outgoing edge and the graph
is acyclic. Choose the options with no incoming edge, tied above the rejected ones, which are tied
among themselves. The conditions use only within-row differences, giving \ax{R}; edges depend only on
principal $2\times2$ blocks and persist under restriction, giving \ax{C}; the rule is neutral, ordered,
and responsive (it chooses $\{1\}$ from $\bigl(\begin{smallmatrix}1&0\\1&0\end{smallmatrix}\bigr)$ and
both options from the zero matrix); and raising a diagonal entry only removes incoming and adds
outgoing edges, so it satisfies \Mon. It fails \ax{F}, since scaling by a small $a>0$ removes every
edge.

With \ax{F}, given $M\in[0,1]^{n\times n}$ and small $\alpha>0$, pass from $M$ to $\alpha M$ by \ax{F} and
then to $A_{ij}=\tfrac12+\alpha(M_{ij}-M_{ii})$ by \ax{R}; every intermediate matrix stays in $[0,1]$.
Embed $A$ in a circulant as in Lemma~\ref{lem:embed}, filling unused residues with $\tfrac12$, and run
the proof of Theorem~\ref{thm:circulant}.
\end{proof}

\begin{proposition}[Prediction bounded away from accuracy]\label{prop:noisy}
Let $P_n=(1-\eta)I+\tfrac{\eta}{n}\one\one^\top$ with $0<\eta<1$, a prediction channel whose
probability of a correct forecast is $1-\eta+\eta/n\le1-\eta/2$ for $n\ge2$. With unrestricted real
payoffs, every $M\in\Rl^{n\times n}$ is the causal matrix of a problem whose predictor uses $P_n$. With
payoffs in $[0,1]$ and \ax{F}, the circulant argument survives, uniformly in $n$.
\end{proposition}

\begin{proof}
$P_n$ is invertible, with $P_n^{-1}=\frac1{1-\eta}I-\frac{\eta}{(1-\eta)n}\one\one^\top$. Let the state be
the forecast label, drawn from row $i$ of $P_n$ when $i$ is chosen, and let $U=P_n^{-1}M$ be the
state-by-action payoff table; then the causal matrix is $P_nU=M$. Since $P_n^{-1}\one\one^\top=\one\one^\top$,
the constant matrix $\tfrac12\one\one^\top$ is mapped to itself, and the absolute row sums of $P_n^{-1}$ are
at most $(1+\eta)/(1-\eta)$; so a neighbourhood of the constant matrix whose radius does not depend on
$n$ is mapped into $[0,1]$, and the compressed matrices of Proposition~\ref{prop:bounded} remain
realisable.
\end{proof}

An accuracy bound alone, then, does not destroy richness. Completely uninformative prediction is
different: if all induced laws coincide, the channel loses rank, every row of the causal matrix is
the same, and the problem is an ordinary causal one.

\begin{proposition}[Bounded action sets]\label{prop:binarytree}
For every $k\ge2$, the symmetric extensions required by Theorem~\ref{thm:circulant} are realisable
by problems with perfect recall and at most $k$ actions per information set, provided contraction
applies to arbitrary sets of policies.
\end{proposition}

\begin{proof}
Take a complete binary decision tree of depth $m$, with the agent observing its own past actions
and the prediction hidden until the end. It has $2^m-1$ information sets and $N=2^m$ action paths,
and each path is produced by exactly $L=2^{2^m-1-m}$ policies, since prescriptions off the path are
free. Let a predictor forecast the eventual path and let the payoff depend on the forecast and the
actual path; any path matrix $C$ is then realisable, and the policy matrix is $C\otimes J_L$, with $J_L$
the all-ones $L\times L$ matrix. If $C$ is circulant, rotating path blocks and permuting duplicate
policies within blocks act transitively on the policies. Any target matrix embeds in a circulant
whose size is a power of two (Lemma~\ref{lem:embed}), and is recovered by restricting to one
representative policy from each designated block.
\end{proof}

The proviso matters: arbitrary sets of surviving policies are generally not obtained by pruning local
action sets, which produces product-shaped sets. Restricting contraction to literal pruning is a
different axiom (Section~\ref{sec:menus}).

\begin{proposition}[A single bounded choice]\label{prop:ceiling}
With one information set and an unbounded action set, every square matrix is realisable. With one
information set and at most $k\ge2$ actions, the rule that is universally indifferent below size $k$
and ranks by $M_{ii}-\frac1k\sum_jM_{ij}$ at size $k$ satisfies \ax{N}, \ax{F}, \ax{R}, \ax{C}, \ax{O},
\ax{Q}, and \Mon.
\end{proposition}

\begin{proof}
The first claim is Theorem~\ref{thm:realise}. For the second, every proper restriction of a
size-$k$ matrix lands where every option is chosen, so \ax{C} holds; the score is a within-row
difference, invariant under \ax{F} up to scale, neutral, responsive, and increasing in $M_{ii}$.
\end{proof}

No finite size ceiling, therefore, forces triviality under the existential form of \ax{Q}. If
non-triviality is required specifically on binary menus, the three-option circulant
$\bigl(\begin{smallmatrix}0&x&y\\y&0&x\\x&y&0\end{smallmatrix}\bigr)$ already blocks it as soon as
three options are available.
